\documentclass[11pt]{article}

\usepackage[a4paper,margin=2.5cm]{geometry}
\usepackage{amsmath,amssymb,amsthm}
\usepackage[colorlinks=true,linkcolor=blue,citecolor=blue,urlcolor=blue]{hyperref}
\usepackage{booktabs}
\usepackage{enumitem}
\usepackage{natbib}

\newtheorem{theorem}{Theorem}[section]
\newtheorem{lemma}{Lemma}[section]
\newtheorem{proposition}{Proposition}[section]
\newtheorem{corollary}{Corollary}[section]
\newtheorem{definition}{Definition}[section]
\newtheorem{remark}{Remark}[section]

\newcommand{\twoI}{2I}
\newcommand{\Vsixhundred}{V_{600}}
\newcommand{\Dicfive}{\mathrm{Dic}_5}
\newcommand{\Vtwentyfour}{V_{24}}
\newcommand{\Ttau}{T_\tau}
\newcommand{\sigmaG}{\sigma}

\title{A $1/4$ Incidence Identity on the 600-Cell\\ and its Bekenstein--Hawking Analogy}
\author{%
  Lee Smart \\[3pt]
  \textit{Institute of Vibrational Field Dynamics} \\[2pt]
  \texttt{contact@vibrationalfielddynamics.org}}
\date{May 2026}

\begin{document}

\maketitle

\noindent\textbf{Status:} Preprint (not peer-reviewed).

\begin{abstract}
We prove a finite-arithmetic identity on the binary icosahedral group
$\twoI$ (equivalently, the 120 vertices of the 600-cell $\Vsixhundred$).
Let $\Dicfive \subset \twoI$ denote the binary dihedral subgroup of order
$20$ generated by the bulk $\Ttau$-cycles, and let $\Vtwentyfour \subset
\Vsixhundred$ denote the 24-cell sublattice fixed by the Galois twist
$\sigmaG : \sqrt{5} \mapsto -\sqrt{5}$ acting on the icosian
quaternion components. We show that
$|gH \cap \Vtwentyfour| = 4$ and $|Hg \cap \Vtwentyfour| = 4$ for every
left coset $gH$ and right coset $Hg$ of $H = \Dicfive$ in $\twoI$, where
$\Dicfive$ is self-normalizing (and hence non-normal) in $\twoI$. Hence
each coset partitions into 4~$\sigmaG$-fixed and 16~$\sigmaG$-mobile
vertices, giving the rational identity
\[
  \frac{S(C)}{A(C)} \;=\; \frac{4}{16} \;=\; \frac{1}{4}
\]
for every coset $C$, where $S(C) := |C \cap \Vtwentyfour|$ and
$A(C) := |C \setminus \Vtwentyfour|$. The identity is moreover
$\tau$-independent: every choice of order-10 element $\tau \in \twoI$
yields a bulk subgroup $H_\tau \cong \Dicfive$, and all such choices
give the same per-coset incidence. Interpreting $S$ as a count of
twin-anchored information bits and $A$ as a count of dynamical horizon
channels, the structural ratio has the same dimensionless numerical
value as the Bekenstein--Hawking entropy coefficient. The result is
parameter-free, established by exact $\mathbb{Q}(\sqrt{5})$ enumeration,
and is independent of any cosmological, dynamical, or quantum-gravity
hypothesis. We make no claim regarding cascade-unit calibration to
physical Planck units; the rational ratio is the load-bearing content.
\end{abstract}

\section{Introduction}\label{sec:intro}

The Bekenstein--Hawking formula, in natural units
$\hbar = G = c = k_B = 1$,
\(
  S_{\mathrm{BH}} = A/4,
\)
\citep{Bekenstein1973,Hawking1975} attaches a thermodynamic entropy
to a black-hole horizon proportional to its area, with the
dimensionless constant of proportionality $1/4$. The coefficient $1/4$ is universal
across asymptotic black-hole spacetimes, and its microscopic origin has
been a central question of quantum gravity.

Two prominent microscopic derivations have been proposed: brane
state-counting in string theory \citep{StromingerVafa1996} and
isolated-horizon spin-network counting in loop quantum gravity
\citep{AshtekarBaezKrasnov2000}. Each yields the $1/4$ coefficient
(modulo standard normalisation conventions in the LQG case) via a
partition function or character formula that involves nontrivial
dynamical machinery.

This paper reports a finite-arithmetic identity that produces the same
dimensionless numerical value by counting alone. Specifically, the count
$|gH \cap \Vtwentyfour|$ equals $4$ for every coset of a particular
order-$20$ subgroup $H$ of the binary icosahedral group, and the
complementary count $|gH \setminus \Vtwentyfour|$ equals $16$. The
rational ratio $4/16 = 1/4$ is therefore an exact structural identity
on a finite group. We do \emph{not} claim a derivation of physical
black-hole thermodynamics; we report a structural ratio that matches
the universal coefficient under a transparent interpretive
identification of the two coset components.

The remainder of the paper is organised as follows.
Section~\ref{sec:model} fixes notation and the exact arithmetic in
$\mathbb{Q}(\sqrt{5})$ used throughout.
Section~\ref{sec:bulk} identifies the bulk subgroup $\Dicfive \subset \twoI$
and gives an explicit binary-dihedral presentation certificate.
Section~\ref{sec:24cell} verifies the $\sigmaG$-fixed sublattice
$\Vtwentyfour$.
Section~\ref{sec:incidence} states and proves the main coset-incidence
theorem. Section~\ref{sec:coefficient} reports the $1/4$ structural
coefficient. Sections~\ref{sec:scope} and~\ref{sec:comparison} mark the
scope and place the result alongside existing microscopic derivations.

\section{Exact model}\label{sec:model}

All arithmetic in this paper is exact. Coordinates of $\Vsixhundred$
vertices lie in the field $K := \mathbb{Q}(\sqrt{5})$, represented as
ordered pairs $(a,b) \in \mathbb{Q}^2$ standing for $a + b\sqrt{5}$.
The Galois automorphism $\sigmaG : K \to K$ acts by
$\sigmaG(a + b\sqrt{5}) = a - b\sqrt{5}$.

\begin{definition}[Icosian quaternions and $\Vsixhundred$]\label{def:icosian}
Let $\mathbb{H}_K$ be the Hamilton quaternion algebra over $K$ and
$\varphi = (1+\sqrt{5})/2$. The \emph{$600$-cell} $\Vsixhundred$ is the
set of $120$ unit icosians obtained as the union of the following three
coordinate orbits with respect to the standard $\mathbb{H}_K$ basis:
\begin{enumerate}[label=(\roman*)]
\item the $8$ \emph{axis} vertices $(\pm 1, 0, 0, 0)$ and the $3$
permutations thereof;
\item the $16$ \emph{half-type} vertices $(\pm 1/2,\, \pm 1/2,\, \pm 1/2,\, \pm 1/2)$;
\item the $96$ \emph{even-permutation} vertices: vectors whose four
coordinates are an even permutation of
$(0,\, \pm 1/(2\varphi),\, \pm 1/2,\, \pm \varphi/2)$ with at least one
zero, and the three nonzero coordinates ranging over all $2^3 = 8$
sign combinations.
\end{enumerate}
This is the standard construction of the binary icosahedral group
$\twoI \subset \mathrm{Sp}(1, K)$ \citep{ConwaySloane1988,Coxeter1973,MoodyPatera1993}.
\end{definition}

The map $\sigmaG$ extends componentwise to icosian quaternions; it
sends $\Vsixhundred$ bijectively to a second 600-cell $\sigmaG(\Vsixhundred)$
that meets $\Vsixhundred$ exactly along the 24 axis- and half-type
vertices identified in Section~\ref{sec:24cell}.

\begin{lemma}[120 vertices, group closure, unit norm]\label{lem:basics}
The set $\Vsixhundred$ has exactly $120$ elements; every element has
$K$-norm $1$; and $\Vsixhundred$ is closed under quaternionic
multiplication, forming the binary icosahedral group $\twoI$.
\end{lemma}

\begin{proof}
Direct enumeration of axis ($8$), half-type ($16$), and even-permutation
($96$) coordinate orbits gives $120$ exact icosian vectors in
$\mathbb{Q}(\sqrt{5})^4$; deduplication confirms cardinality. Unit norm
$|v|^2 = 1$ and exhaustive group closure across all $14{,}400$ products
$v_i v_j$ are checked by the accompanying \texttt{verify.py} certificate
(cf.~\citealp{ConwaySloane1988}).
\end{proof}

We use $G := \twoI = \Vsixhundred$ and $|G| = 120$ throughout.

\section{The bulk subgroup}\label{sec:bulk}

\begin{definition}[$T_\tau$-cycles]\label{def:Ttau}
For $\tau \in G$ with $\operatorname{ord}(\tau) = 10$ (there are
exactly $24$ such elements; pick any one), left multiplication by
$\tau$ generates a permutation $\Ttau : G \to G$, $g \mapsto \tau g$,
whose cycle decomposition partitions $G$ into 12 cycles each of length
10. We will later show (Theorem~\ref{thm:taunoindep}) that the
eventual incidence identity is independent of the choice of $\tau$.
\end{definition}

\begin{definition}[K-classes]\label{def:Kclass}
For each vertex $v \in G$, write $\mathrm{shell}(v)$ for the index
($0$ through $8$) of the rational squared-distance class
$|v - 1|^2 \in K$ in the totally ordered set of distinct shell radii;
ordering is by exact sign comparison on $K$ (no floating point).
For a cycle $C$, set
\(
  K(C) := \sum_{v \in C}(\mathrm{shell}(v) - 4)^2.
\)
\end{definition}

\begin{lemma}[K-multiset]\label{lem:kmulti}
The 12 cycles distribute into the K-multiset
$\{72{:}1,\, 0{:}1,\, 52{:}5,\, 20{:}5\}$: one cycle with $K = 72$
(the maximum), one with $K = 0$ (vertices uniformly on shell $4$),
five with $K = 52$, and five with $K = 20$.
\end{lemma}

\begin{proof}
Direct rational computation on $G$.
\end{proof}

\begin{definition}[Bulk subgroup $H$]\label{def:bulk}
Let $H \subset G$ be the union of the two K-special cycles ($K = 72$ and
$K = 0$). Then $|H| = 20$.
\end{definition}

\begin{lemma}[$H = \Dicfive$]\label{lem:h_is_dic5}
$H$ is closed under quaternionic multiplication, contains the identity,
and admits the binary-dihedral presentation
\(
  H \;=\; \langle a, b \;|\; a^{10} = 1,\ b^2 = a^5,\ b a b^{-1} = a^{-1} \rangle
  \;\cong\; \Dicfive.
\)
\end{lemma}

\begin{proof}
Closure and presence of $1 \in H$ are verified by exhaustive enumeration
across all $400$ products in $H \times H$. The generator search returns
an explicit pair $(a, b) \in H \times H$: $a$ is an order-$10$ element
drawn from the $K = 72$ cycle, and $b$ is one of ten elements of $H$
satisfying $b^2 = a^5$ and $bab^{-1} = a^{-1}$ (we fix one such choice).
The accompanying certificate then verifies $\langle a, b \rangle = H$
by enumerating the closure of $\{a, b\}$ under quaternionic
multiplication and comparing it to $H$.
\end{proof}

\begin{lemma}[Coset cycle composition]\label{lem:cosetkpair}
For $T_\tau : g \mapsto \tau g$ with $\tau \in H$, every non-bulk
right coset $Hg$ is the disjoint union of one whole $T_\tau$-cycle of
K-class $52$ and one whole $T_\tau$-cycle of K-class $20$. The five
non-bulk right cosets therefore exhaust the ten non-bulk
$T_\tau$-cycles. By contrast, every non-bulk left coset $gH$ contains
exactly two vertices from each of the ten non-bulk cycles, hence
$10$ vertices of K-class $52$ and $10$ of K-class $20$ in aggregate
but is not a union of whole $T_\tau$-cycles.
\end{lemma}

\begin{proof}
Right cosets: $T_\tau(Hg) = (\tau H)g = Hg$ because $\tau \in H$
implies $\tau H = H$ (closure of $H$). Therefore $Hg$ is closed under
$T_\tau$ and decomposes as a disjoint union of whole $T_\tau$-cycles;
the cardinality count $|Hg| = 20 = 10 + 10$ together with the K-class
assignment gives one cycle of each non-bulk K-class.

Left cosets: $T_\tau(gH) = \tau g H$, which equals $gH$ iff
$g^{-1}\tau g \in H$. The certificate \texttt{verify.py} checks the
finite condition $g^{-1}\tau g \notin H$ for every $g \in G \setminus H$
explicitly. This direct enumeration (not derivable from
self-normality of $H$ alone) establishes that $T_\tau$ does not
preserve any non-bulk left coset. The induced action of $T_\tau$ on
the five non-bulk left cosets is therefore a permutation with no
fixed points; since $\tau^{10} = 1$ acts trivially on cosets and the
order divides $10$, the induced permutation has order $5$, i.e.\ it
is a $5$-cycle. Hence the $T_\tau$-orbit of any vertex
$v \in gH$ has length $10$ and visits each of the five non-bulk left
cosets exactly twice, giving the $2$-per-cycle distribution. The
aggregate K-class composition $10$+$10$ in each left coset follows
immediately.

The whole-cycle / 2-per-cycle distinction is verified by the shared
tests
\texttt{test\_non\_bulk\_left\_cosets\_2\_per\_cycle} and
\texttt{test\_non\_bulk\_right\_cosets\_pair\_K52\_K20}.
\end{proof}

\begin{lemma}[$H$ is self-normalizing]\label{lem:nonnormal}
For every $g \in G \setminus H$, the left coset $gH$ and the right coset
$Hg$ are distinct subsets of $G$. Hence $N_G(H) = H$, and $H$ is
self-normalizing and \emph{not} normal in $G$.
\end{lemma}

\begin{proof}
Direct enumeration: there are $|G| - |H| = 100$ elements outside $H$,
and the certificate \texttt{verify.py} confirms $gH \neq Hg$ for each
of them. Therefore the only $g$ with $gH = Hg$ are those in $H$
itself, giving $N_G(H) = H$. Self-normalizing subgroups are non-normal
unless they coincide with $G$, which $H$ does not.
\end{proof}

A consequence of Lemma~\ref{lem:nonnormal} is that the main incidence
result of Section~\ref{sec:incidence} cannot be inferred from a quotient
construction: it must be checked on each coset family separately. We do
so explicitly.

\section{The $\sigmaG$-fixed 24-cell}\label{sec:24cell}

\begin{definition}[Galois twist on icosians]\label{def:galois}
$\sigmaG$ acts componentwise on icosian quaternions: for
$v = (v_w, v_x, v_y, v_z)$ with each $v_i \in K$,
$\sigmaG(v) := (\sigmaG(v_w), \sigmaG(v_x), \sigmaG(v_y), \sigmaG(v_z))$.
\end{definition}

\begin{definition}[24-cell as $\sigmaG$-fixed sublattice]\label{def:V24}
$\Vtwentyfour := \{\, v \in \Vsixhundred \;:\; \sigmaG(v) = v \,\}$.
\end{definition}

A vertex $v \in \Vsixhundred$ lies in $\Vtwentyfour$ if and only if
every coordinate has zero $\sqrt{5}$-component, i.e.\ if every
coordinate is rational.

\begin{lemma}[$|\Vtwentyfour| = 24$]\label{lem:24cell}
$\Vtwentyfour$ has cardinality $24$ and consists of the $8$ axis points
$(\pm 1, 0, 0, 0)$ (and permutations) and the $16$ half-type points
$(\pm 1/2, \pm 1/2, \pm 1/2, \pm 1/2)$. These $24$ points form the
exceptional regular $24$-cell, the only regular convex 4-polytope
without a 3-dimensional analogue \citep{Coxeter1973}.
\end{lemma}

\begin{proof}
The axis points are coordinatewise rational. The half-type points have
coordinates $\pm 1/2$. The 96 even-permutation points have at least one
coordinate involving $\varphi/2 = (1 + \sqrt{5})/4$ or
$1/(2\varphi) = (\sqrt{5} - 1)/4$, hence a non-zero $\sqrt{5}$
component. Direct enumeration confirms $8 + 16 = 24$.
\end{proof}

\section{The incidence theorem}\label{sec:incidence}

\begin{theorem}[Coset incidence]\label{thm:main}
Let $G = \twoI$ and $H = \Dicfive \subset G$ as in
Lemma~\ref{lem:h_is_dic5}, and let $\Vtwentyfour \subset G$ be the
$\sigmaG$-fixed $24$-cell of Lemma~\ref{lem:24cell}. Then for every
left coset $gH$ and every right coset $Hg$ of $H$ in $G$,
\[
  |gH \cap \Vtwentyfour| \;=\; |Hg \cap \Vtwentyfour| \;=\; 4,
\]
and consequently
\[
  |gH \setminus \Vtwentyfour| \;=\; |Hg \setminus \Vtwentyfour| \;=\; 16.
\]
\end{theorem}

\begin{proof}
$|H| = 20$ (Lemma~\ref{lem:h_is_dic5}), so each coset has cardinality
$20$ by Lagrange. There are $|G|/|H| = 6$ left cosets and $6$ right
cosets, partitioning $G$ into $6 \times 20 = 120$ vertices.

The intersection $H \cap \Vtwentyfour$ has exactly $4$ elements
(direct enumeration: the identity $1$, the central element of order $2$,
and two elements of order $4$, all coordinatewise rational). Let $C$
be a non-bulk left coset; we show $|C \cap \Vtwentyfour| = 4$ by exact
$\mathbb{Q}(\sqrt{5})$ enumeration of $C$ and explicit intersection
with $\Vtwentyfour$. The verification script \texttt{verify.py}
performs this enumeration over both coset families, exhibits the
explicit four-element intersection by vertex index in each case, and
checks that the six left cosets (resp.\ six right cosets) partition
$G$. The result is summarised in Table~\ref{tab:incidence}.

Because $H$ is non-normal (Lemma~\ref{lem:nonnormal}), the right-coset
incidence is verified independently. The conclusion
$|Hg \cap \Vtwentyfour| = 4$ is established by the analogous explicit
enumeration. The complementary count
$|gH \setminus \Vtwentyfour| = |gH| - 4 = 16$ follows immediately, and
similarly for right cosets.
\end{proof}

\begin{theorem}[$\tau$-independence]\label{thm:taunoindep}
The conclusion of Theorem~\ref{thm:main} is independent of the choice
of order-$10$ element $\tau$ used to construct the bulk subgroup $H$.
Concretely, for each of the $24$ order-$10$ elements $\tau \in G$,
the resulting bulk subgroup $H_\tau := (K_\tau = 72) \cup (K_\tau = 0)$
is binary-dihedral of order $20$, and every left and right coset of
$H_\tau$ in $G$ contains exactly $4$ vertices of $\Vtwentyfour$. The
$24$ choices yield exactly $6$ distinct subgroups $H_\tau$, all of
which are conjugate in $G$.
\end{theorem}

\begin{proof}
For each of the $24$ order-$10$ elements, the verification script
constructs $H_\tau$, partitions $G$ into the six $H_\tau$-cosets, and
checks the 4-fixed/16-mobile incidence with respect to $\Vtwentyfour$.
The enumeration produces six distinct $H_\tau$, all order-$20$ binary
dihedral, all $G$-conjugate (by inspection of the orbit of $H$ under
conjugation by elements of $G$). The 4/16 incidence is uniform across
all $24$ choices.
\end{proof}

\begin{table}[h]
\centering
\caption{Coset incidence on $G = \twoI$ with $H = \Dicfive$ and
$\Vtwentyfour$ the $\sigmaG$-fixed sublattice. The incidence is uniform
across left and right coset families; $H$ is non-normal so this is not
a quotient identity.}
\label{tab:incidence}
\begin{tabular}{c r r r r r}
\toprule
side & coset & $|C|$ & $|C \cap \Vtwentyfour|$ & $|C \setminus \Vtwentyfour|$ & $S/A$ \\
\midrule
left  & $0\;(=H)$ & 20 & 4 & 16 & 1/4 \\
left  & 1 & 20 & 4 & 16 & 1/4 \\
left  & 2 & 20 & 4 & 16 & 1/4 \\
left  & 3 & 20 & 4 & 16 & 1/4 \\
left  & 4 & 20 & 4 & 16 & 1/4 \\
left  & 5 & 20 & 4 & 16 & 1/4 \\
\midrule
right & $0\;(=H)$ & 20 & 4 & 16 & 1/4 \\
right & 1 & 20 & 4 & 16 & 1/4 \\
right & 2 & 20 & 4 & 16 & 1/4 \\
right & 3 & 20 & 4 & 16 & 1/4 \\
right & 4 & 20 & 4 & 16 & 1/4 \\
right & 5 & 20 & 4 & 16 & 1/4 \\
\bottomrule
\end{tabular}
\end{table}

\section{The structural coefficient}\label{sec:coefficient}

\begin{definition}[Anchor and channel counts]\label{def:SAcount}
For a subset $C \subseteq G$, set
\(
  S(C) := |C \cap \Vtwentyfour|, \quad
  A(C) := |C \setminus \Vtwentyfour|.
\)
We refer to $S(C)$ as the \emph{anchor count} and $A(C)$ as the
\emph{channel count} of $C$.
\end{definition}

\begin{theorem}[$1/4$ identity]\label{thm:onequarter}
For every coset $C$ of $H$ (left or right) in $G$,
\(
  \dfrac{S(C)}{A(C)} = \dfrac{1}{4}.
\)
For the union $B$ of the five non-bulk left cosets,
$S(B) = 20$, $A(B) = 80$, and $S(B)/A(B) = 1/4$.
\end{theorem}

\begin{proof}
Combine Theorem~\ref{thm:main} (numerator and denominator $4$ and $16$
per coset) with the partition of $B$ into five disjoint cosets
($S$ and $A$ are additive over disjoint unions).
\end{proof}

\section{Scope}\label{sec:scope}

\noindent
\fbox{\parbox{\dimexpr\linewidth-2\fboxsep-2\fboxrule\relax}{%
\textbf{Scope.}
\begin{itemize}
\item \textbf{Established (exact, unconditional):} the rational identity
$|gH \cap \Vtwentyfour| / |gH \setminus \Vtwentyfour| = 1/4$ on the
finite group $G = \twoI$ with $H = \Dicfive$, for both left and
right cosets, with $H$ self-normalizing and non-normal; and the
$\tau$-independence of the identity over all $24$ order-$10$ choices.
\item \textbf{Interpretive:} the identification of $S$ with anchor
information bits and $A$ with horizon channels. This is offered as
a definition with motivation, not as a theorem. A different
interpretation may yield a different physical content from the same
arithmetic identity.
\item \textbf{Out of scope:} cascade-unit calibration to physical
Planck units; the canonical $\tau_\sigma$ involution (assigned to a
companion paper); derivation of the Hawking spectrum;
Friedmann-equation analog; cosmological observables; ontology of
$\Vsixhundred$. The numerical match between the structural ratio
$1/4$ and the Bekenstein--Hawking coefficient $1/4$ is observed
without claim; the unit-calibration mapping required to lift this
from a numerical coincidence to a quantitative physical derivation
is not constructed here.
\end{itemize}}}

\section{Comparison with existing microscopic derivations}\label{sec:comparison}

We do not claim superiority over the established microscopic
derivations; we report a methodologically distinct route. The
following table summarises the relationship.

\begin{table}[h]
\centering
\begin{tabular}{l l l}
\toprule
Approach & Mathematical content & Output (in natural units, $\hbar = G = c = k_B = 1$) \\
\midrule
Strominger--Vafa \citep{StromingerVafa1996} & D-brane partition function & $S/A = 1/4$ \\
Ashtekar--Baez--Krasnov \citep{AshtekarBaezKrasnov2000} & spin-network state count & $S/A = 1/4$ (after standard LQG normalisation) \\
This paper & finite-group incidence on $\twoI$ & $S/A = 1/4$ (dimensionless) \\
\bottomrule
\end{tabular}
\end{table}

The incidence approach computes a dimensionless ratio. To match
physical units, one would need to identify the cascade ``area'' unit
(one $\sigmaG$-mobile vertex) with a Planck-area count and the
cascade ``entropy'' unit (one $\sigmaG$-fixed vertex) with one bit of
horizon information. The calibration constant is not derived here; it
is the analogue of the $\hbar G c^{-3}$ scaling in the standard
formula. What is established unconditionally is that the structural
ratio is rational and equal to $1/4$.

\section{Conclusion}\label{sec:concl}

The Bekenstein--Hawking coefficient $1/4$ has a finite-arithmetic
analogue as a coset-incidence identity on the binary icosahedral
group $\twoI$, namely $|gH \cap \Vtwentyfour|/|gH \setminus \Vtwentyfour|
= 4/16 = 1/4$ for every coset of the binary dihedral subgroup
$H = \Dicfive$ against the $\sigmaG$-fixed $24$-cell sublattice
$\Vtwentyfour$, with both left and right coset families behaving
identically despite the non-normality of $H$, and with the identity
$\tau$-independent across all $24$ admissible bulk constructions.
The result is exact ($\mathbb{Q}(\sqrt{5})$ arithmetic, no floating
point), parameter-free, and verifiable by the accompanying script. We
state explicitly that this is a structural identity rather than a
derivation of physical black-hole thermodynamics, and we leave the
cascade-unit calibration to subsequent work.

\appendix

\section{Verification certificate}\label{app:verify}

The companion script \texttt{verify.py} performs the following exact
checks and exits with status $0$ when all assertions hold:
\begin{enumerate}
\item $|\Vsixhundred| = 120$.
\item Unit norm $|v|^2 = 1$ for all $120$ vertices.
\item \textbf{Exhaustive} group closure: every product $v_i v_j$ across
all $120^2 = 14{,}400$ pairs lies in $\Vsixhundred$.
\item $T_\tau$-cycle decomposition: 12 cycles of length 10; K-multiset
$\{72{:}1, 0{:}1, 52{:}5, 20{:}5\}$.
\item $|H| = 20$, $1 \in H$, $H$ closed under multiplication (all $400$
$H \times H$ products), $H$ self-normalizing in $G$.
\item Explicit Dic$_5$ generators $(a, b)$ with $a \in H$ of order
$10$, $b \in H$ satisfying $b^2 = a^5$ and $bab^{-1} = a^{-1}$, and
$\langle a, b \rangle = H$ verified by closure enumeration.
\item $|\Vtwentyfour| = 24$ and $\Vtwentyfour = \Vsixhundred \cap
\sigma(\Vsixhundred)$, with $|H \cap \Vtwentyfour| = 4$ and
$|\Vtwentyfour \setminus H| = 20$.
\item Per-coset incidence table for both left and right cosets
(Table~\ref{tab:incidence}), with explicit $4$-element intersection
indices printed for each coset.
\item Aggregate boundary $S/A = 20/80 = 1/4$.
\item $\tau$-independence sweep over all $24$ order-$10$ elements
$\tau \in G$: each yields a binary-dihedral $H_\tau$ of order $20$
with the explicit $(a, b)$ presentation, partitions $G$ into six
left and six right cosets of size $20$ each with uniform
$4 / 16$ incidence, and the $24$ choices yield exactly $6$ distinct
$H_\tau$ subgroups all $G$-conjugate to the canonical $H$.
\end{enumerate}

The script depends only on the Python standard library
(\texttt{fractions}, \texttt{itertools}, \texttt{collections}) and on
the shared \texttt{vfd\_v600} package
(\texttt{papers/v600-programme/lib}).

Reproducibility command (from the repository root):
\begin{verbatim}
PYTHONPATH=papers/v600-programme/lib \
    python3 papers/bekenstein-incidence/verify.py
\end{verbatim}

\bibliographystyle{plainnat}
\bibliography{references}

\end{document}
